% !TeX root = ../../progress-sem1.tex
% Section 3.4 (subsections): one subsection per layer of the system architecture,
% matching "สถาปัตยกรรมโดยภาพรวมของระบบ" in Progress_Sem1_THAI_Example1.pdf.
% Input from chapters/03-methodology.tex, after the architecture figure.

\subsection{Client Side (ส่วนผู้ใช้งาน)}
\thaijustify{
  ผู้ใช้ทุกบทบาทเข้าใช้งานผ่านเว็บเบราว์เซอร์ นักออกแบบ (Designer) ใช้หน้าจอหลักของแพลตฟอร์มตั้งแต่การสร้างโปรเจกต์จนถึงการปรับแก้ UI ผู้เข้าร่วมทดสอบ (Usability Tester) เข้าร่วมเซสชันผ่านลิงก์เชิญโดยไม่ต้องสมัครสมาชิก ส่วน Client, Developer และ Project Manager ใช้มุมมองสำหรับติดตามสถานะ รับมอบงาน และดูรายงาน
}

\subsection{Frontend Layer}
\thaijustify{
  Frontend พัฒนาด้วย Next.js แบบ App Router แบ่ง route ออกเป็นกลุ่ม \texttt{(auth)} สำหรับการเข้าสู่ระบบ \texttt{(dashboard)} สำหรับ workspace และโปรเจกต์ และ \texttt{(legal)} สำหรับนโยบายความเป็นส่วนตัวและข้อกำหนดการใช้งาน หน้าจอสื่อสารกับ backend ผ่าน REST API และรับความคืบหน้าของงาน AI ผ่าน SSE หน้าตัวอย่าง UI ที่ AI สร้างแสดงใน iframe ที่ไม่อนุญาตให้รันสคริปต์ เพื่อป้องกันโค้ดที่ไม่ปลอดภัยจากผลลัพธ์ของโมเดล
}

\subsection{Backend API Layer}
\thaijustify{
  Backend พัฒนาด้วย NestJS แบบ monolith แบ่งเป็น module ตามโดเมน ได้แก่ \texttt{auth}, \texttt{access}, \texttt{users}, \texttt{workspaces}, \texttt{projects}, \texttt{design}, \texttt{rounds}, \texttt{generation}, \texttt{ai}, \texttt{testing}, \texttt{iteration}, \texttt{stakeholders}, \texttt{audit} และ \texttt{health} ทุก request ผ่านการตรวจ token การตรวจสิทธิ์ตาม permission การจำกัดอัตราการเรียก (rate limit) และ security header ก่อนถึง module ปลายทาง
}

\subsection{AI Job Processor Layer}
\thaijustify{
  งานที่ใช้เวลานานถูกส่งเข้าคิวของ BullMQ บน Redis แทนการประมวลผลใน request โดยตรง โดยแบ่งเป็นสองคิว ได้แก่ คิว design-generation สำหรับการสร้าง UI variant และการสังเคราะห์ persona และคิว ai-task สำหรับการร่าง user story และการวิเคราะห์ข้อเสนอแนะ worker ดึงงานจากคิว เรียกโมเดลผ่าน CPE AI Gateway ซึ่งมี API แบบเดียวกับ OpenAI และส่งความคืบหน้ากลับไปยังหน้าจอผ่าน SSE ผู้ใช้จึงปิดหน้าจอแล้วกลับมาติดตามงานเดิมต่อได้ รายละเอียดอยู่ในหัวข้อ~\ref{sec:ai-pipeline}
}

\subsection{Data \& External Services Layer}
\thaijustify{
  PostgreSQL จัดเก็บข้อมูลทั้งหมดของระบบผ่าน Prisma (หัวข้อ~\ref{sec:er-diagram}) Redis ทำหน้าที่ทั้งคิวงานและ cache ของ response บริการภายนอกที่ใช้ ได้แก่ Google OAuth 2.0 สำหรับการเข้าสู่ระบบ และ CPE AI Gateway ของภาควิชาซึ่งเข้าถึงได้ผ่าน VPN ของมหาวิทยาลัยเท่านั้น
}
